\section{\texorpdfstring{哈密顿量 $H$ 的边界深度}{哈密顿量 H 的边界深度}}
\label{sec:hamiltonian-boundary-depth}

在着手构造谱不变量 $\rho(H;a)$ 之前，先定义哈密顿量 $H$ 所伴随的一个更为朴素的
\stfhIndex{I043}{00-beta-h}{$\beta(H)$}
不变量，即\termBoundaryDepth{}~$\beta(H)$；它是 13.8 节定义 13.8.5 的特殊情形。
我们将证明，$\beta(H)$ 是滤过链复形 $CF_*(H)$ 的同构型不变量。

\begin{definition}[边界深度]
\label{def:21-hamiltonian-boundary-depth}
设 $H$ 是给定的非退化哈密顿量，并设 $(H,J)$ 弗洛尔正则。对相应弗洛尔复形
$(CF_*(H,J),\partial_{(H,J)})$，定义
$\beta(H,J):=\beta(\partial_{(H,J)})$，即
\[
 \beta(H,J)=\inf\bigl\{\beta\in\mathbb R\mid
 \forall\lambda\in\mathbb R,\space
 CF_*^\lambda(H,J)\cap\partial CF_*(H,J)
 \subset\partial(CF_*^{\lambda+\beta}(H,J))\bigr\}.
\]
\end{definition}

下述引理立即来自推论~\ref{cor:21-change-almost-complex-structure}。

\begin{lemma}
\label{lem:21-boundary-depth-independent-j}
设 $H$ 非退化，且 $(H,J^0)$、$(H,J^1)$ 均弗洛尔正则，则
\[
 \beta(H,J^0)=\beta(H,J^1).
\]
\end{lemma}

\begin{proof}
证明 $\beta(H,J^0)\leq\beta(H,J^1)$，反向同理。只须证明：若 $\beta\geq0$ 满足
\[
 CF_*^\lambda(H,J^0)\cap\partial CF_*(H,J^0)
 \subset\partial(CF_*^{\lambda+\beta}(H,J^0)),
\]
则对任意 $\epsilon>0$，
\begin{equation}
 CF_*^\lambda(H,J^1)\cap\partial CF_*(H,J^1)
 \subset\partial(CF_*^{\lambda+\beta+\epsilon}(H,J^1)).
 \label{eq:21-boundary-depth-transfer}
\end{equation}
选取从 $J^0$ 到 $J^1$、相对于 $H$ 正则的道路 $j$，记相应链映射为
$h_{(H,j)}:CF_*(H,J^0)\to CF_*(H,J^1)$。由推论
\ref{cor:21-change-almost-complex-structure}，对每个
$\lambda\in\mathbb R\setminus\operatorname{Spec}(H)$，它限制为链同构
\footnote{译者注：原文称 $h_{(H,j)}^\lambda$ 为链同构；但所引推论
\ref{cor:21-change-almost-complex-structure} 只断言它是链映射且在同调上诱导同构。
译文照录原断言。}
\[
 h_{(H,j)}^\lambda:CF_*^\lambda(H,J^0)\longrightarrow CF_*^\lambda(H,J^1),
\]
且
\[
 h_{(H,j)}^\lambda\circ\partial_{(H,J^0)}
 =\partial_{(H,J^1)}\circ h_{(H,j)}^\lambda.
\]
因此，若 $\lambda+\beta\notin\operatorname{Spec}(H)$，则
\eqref{eq:21-boundary-depth-transfer} 甚至可取 $\epsilon=0$；否则可取任意满足
$\lambda+\beta+\epsilon\notin\operatorname{Spec}(H)$ 的 $\epsilon>0$。因
$\mathbb R\setminus\operatorname{Spec}(H)$ 在 $\mathbb R$ 中稠密，$\epsilon$ 可任意小。
交换 $J^0,J^1$ 即完成证明。
\end{proof}

据此，今后简记 $\beta(H)=\beta(H,J)$，其中 $J$ 是使 $(H,J)$ 弗洛尔正则的任一
\stfhIndex{I190}{Eshe}{厄舍（Usher）}
选择。厄舍在 (Ush11) 中证明了边界深度的下列基本性质。（与第 (2) 项相似的结论
见 (Oh09a) 的命题 8.8。）

\begin{theorem}[厄舍]
\label{thm:21-boundary-depth-properties}
边界深度 $\beta$ 满足：
\begin{enumerate}[label=(\arabic*)]
 \item 若 $H,K$ 是两个非退化哈密顿量，则
 \[
  |\beta(H)-\beta(K)|\leq\lVert H-K\rVert.
 \]
 特别地，$\beta$ 关于 $L^{(1,\infty)}$ 拓扑利普希茨连续，因而在该拓扑下连续
 延拓为非负函数。
 \item $\beta(H)\leq\lVert H\rVert$。
 \item 若非退化哈密顿量 $H,K$ 满足 $H\sim K$，则 $\beta(H)=\beta(K)$。
\end{enumerate}
\end{theorem}

\begin{proof}
考虑线性同伦 $\mathcal H^{\mathrm{lin}}:s\mapsto(1-s)H+sK$。可为之选取适当的
$\mathcal J$，使 $(\mathcal H^{\mathrm{lin}},\mathcal J)$ 弗洛尔正则。仍然只须证明
\begin{equation}
 \beta(H)\leq\beta(K)+\lVert H-K\rVert
 \label{eq:21-boundary-depth-lipschitz}
\end{equation}
及交换 $H,K$ 后的版本。沿用引理~\ref{lem:21-boundary-depth-independent-j} 的论证，
这次对 $\mathcal H=\mathcal H^{\mathrm{lin}}$ 使用推论
\ref{cor:21-filtered-continuation-map}，即得~\eqref{eq:21-boundary-depth-lipschitz}，从而
证明第 (1) 项。令 $K\to0$ 得第 (2) 项。最后，引理
\ref{lem:21-boundary-depth-independent-j} 的同一证明适用于第 (3) 项；此时使用
21.6.2 节的推论~\ref{cor:21-optimal-continuation-filtration}。
\footnote{译者注：原文写“21.6.2 节的推论 21.1.5”，但推论 21.1.5 不在该节；
按论证所需的保滤过结论改引推论 21.6.5。}
\end{proof}

正如厄舍 (Ush11) 所指出的，该定理说明 $\beta(H)$ 是滤过链复形 $CF_*(H,J)$ 的
同构型不变量。边界深度与下一节将构造的谱不变量 $\rho(H;a)$ 有一点有趣的区别：
它确实是复形 $CF_*(H,J)$ 的不变量，却不像谱不变量那样出自同调层面的理由。
\footnote{译者注：原文此处写作“将在下一章构造”，但谱不变量紧接着在本章
21.3 节构造；译文按实际结构改作“下一节”。}
关于边界深度在辛拓扑和哈密顿动力系统各类问题中的有趣应用，参见 (Ush11)、
(Ush12)、(Ush13)。
